\chapter{Kinetics of Continuous Media}

Recall the definition of admissible motion
\begin{figure}[H]
\centering
\resizebox{0.8\textwidth}{!}{\input{figure22.eps_tex}}
\caption{Scheme of deformation}
\label{Scheme of deformation}
\end{figure}
\[\hat{\tenq{y}}(\bullet,t)\colon R\to R_t\]
\[t\in\tau=[t_0,t_1]\]
\[\tenq{y}=\hat{\tenq{y}}(\tenq{x},t)=\tenq{x}+\tenq{u}(\tenq{x},t)\]
\[\hat{\tenq{y}}\in C^2(R\times\tau)\]
\[J = \det\nabla\hat{\tenq{y}} = \det\left(\tenq{1}+\nabla\tenq{u}\right)>0\]
\[\hat{\tenq{y}}^{-1}\equiv\hat{\tenq{x}}(\bullet,t)\text{ on }R_t, \hat{\tenq{x}}(\bullet,t)\in C^2(R_t)\]

\section{Mass density, conservation of mass (equation of continuity)}

Assume that exists a scalar valued function $\rho(\bullet,t)$ on $R_t$, $\forall t \in[t_0,t_1]$ called the mass density of the body at time $t$, with the following properties. 

\begin{enumerate}[label=(\roman*)]

\item 
$\rho(\bullet,t)\in C(R_t)$, $\rho(\tenq{y},t)>0\phantom{-} \forall \tenq{y}\in R_t \phantom{-} \forall t\in [t_0,t_1]$
\item \label{conservation of mass-item 2}
If $P\subset R$ is a compact region and if $P_t = \hat{\tenq{y}}(P,t)$, then 

\[m(P_t) = \int_{P_t}\rho(\tenq{y},t)d\volume_{\tenq{y}}\phantom{-} \forall t\in\tau\]

dim $\rho = [ML^{-3}]$
\begin{figure}[H]
\centering
\resizebox{0.6\textwidth}{!}{\input{figure23.eps_tex}}
\caption{Scheme of deformation of $P$}
\label{Scheme of deformation of P}
\end{figure}

\end{enumerate}

\section{Postulate of mass conservation (Global Version)}


\begin{equation}
\boxed{m(P_t) = m(P_{t_0}), t_0\leq t\leq t_1, \forall \text{ compact }P \subset R} \tag{$\ast$} \label{mass conservation-*}
\end{equation}
i.e., $m(P_{t_0})$ is time independent. 

\begin{claim}\mbox{}\\

Define

\begin{equation}
\phi(\tenq{x},t) = \rho(\hat{\tenq{y}}(\tenq{x},t),t)J(\tenq{x},t)\phantom{-} \forall (\tenq{x},t)\in R \times \tau \tag{a} \label{mass conservation-a}
\end{equation}

Then from $\eqref{mass conservation-*}$

\begin{equation}
\phi(\tenq{x},t) = \phi(\tenq{x},t_0)\equiv\footnotemark\rho_0(\tenq{x})\phantom{-}\forall (\tenq{x},t)\in R \times \tau \tag{b} \label{mass conservation-b}
\end{equation}
\footnotetext{no longer depend on time}
Also

\begin{equation}
m(P_t) = \int_P \rho_0(\tenq{x})d\volume_{\tenq{x}}\tag{c} \label{mass conservation-c}
\end{equation}

Finally

\begin{equation}
\rho(\tenq{y}\footnotemark,t)J(\tenq{x},t) = \rho_0(\tenq{x})\phantom{-}\forall (\tenq{x},t)\in R \times \tau \tag{d} \label{mass conservation-d}
\end{equation}
\footnotetext{$\tenq{y} = \hat{\tenq{y}}(\tenq{x},t)$}
 
\end{claim}

\begin{proof}\mbox{}\\

\begin{equation}\tag{1}\label{mass conservation-1}
\begin{split}
\text{\ref{conservation of mass-item 2}} \Rightarrow 
m(P_t)
&=\int_P \rho(\hat{\tenq{y}}(\tenq{x},t),t)J(\tenq{x},t)d\volume_{\tenq{x}}\\
&\vc{=}{\eqref{mass conservation-a}}{}\int_P \phi(\tenq{x},t)d\volume_{\tenq{x}}
\end{split}
\end{equation}

\begin{align*}
\eqref{mass conservation-1}\&\eqref{mass conservation-*}
&\Rightarrow m(P_t) - m(P_{t_0}) = \int_P \left\{\phi(\tenq{x},t)-\phi(\tenq{x},t_0)\right\}d\volume_{\tenq{x}} = 0\\
&\Rightarrow \phi(\tenq{x},t) = \phi(\tenq{x},t_0)\equiv\rho_0(\tenq{x})
\end{align*}

From \ref{conservation of mass-item 2}
\[m(P_t) = \int_{P_t}\rho(\tenq{y},t)d\volume_{\tenq{y}} = \int_P \phi(\tenq{x},t)d\volume_{\tenq{x}} = \int_P \rho_0(\tenq{x})d\volume_{\tenq{x}}\]

\[
\Rightarrow
\boxed{
\begin{aligned}
& \rho(\tenq{y},t)J(\tenq{x},t)=\rho_0(\tenq{x})\\
 \text{ or }&\rho(\hat{\tenq{y}}(\tenq{x},t),t)J(\tenq{x},t)=\rho_0(\tenq{x})
\end{aligned}
}\tag{$\ast\ast$} \label{mass conservation-**}
\]

$\eqref{mass conservation-**}\cdots$ Local version of conservation of mass\\
$\rho_0(\tenq{x})\cdots$ \emph{nominal} or \emph{referential} mass density


\end{proof}

\begin{remark}\mbox{}\\

\begin{enumerate}[label=(\arabic*)]
\item incompressible body ($J(\tenq{x},t)=1$)\\
(can only sustain locally volume preserving deformation)
\[\Rightarrow\rho(\tenq{y},t) = \rho_0(\tenq{x})\]
\item infinitesimal deformation $\abs{\nabla\tenq{u}}\ll 1$, $\tenq{E}\sim \tenq{\varepsilon} = \frac{1}{2}\left(\nabla \tenq{u} + \left(\nabla \tenq{u}\right)^T\right)$

\[\Rightarrow J = \det\left(\tenq{1} + \nabla \tenq{u}\right)\approx 1 + \tr\tenq{\varepsilon} = 1+\cancelto{\text{neglected}}{u_{k,k}}\]
\[\Rightarrow\rho(\tenq{y},t) = \rho_0(\tenq{x})\]
\end{enumerate}
\end{remark}

\section{Velocity and acceleration fields}

The material (Lagrangian, referential) velocity and acceleration defined as 

\begin{figure}[H]
\centering
\resizebox{0.8\textwidth}{!}{\input{figure24.eps_tex}}

\caption{the motion of a body}
\label{Motion of a body-1}
\end{figure}

\[\dot{\tenq{y}}(\tenq{x},t) = \frac{\partial}{\partial t}\hat{\tenq{y}}(\tenq{x},t) = \tenq{v}(\tenq{x},t)\]
\[\ddot{\tenq{y}}(\tenq{x},t) = \frac{\partial^2}{\partial t^2}\hat{\tenq{y}}(\tenq{x},t) = \tenq{a}(\tenq{x},t)\]

where $\tenq{v}(\tenq{x},t)=\dot{\tenq{y}}(\tenq{x},t)\cdots$ velocity at time $t$ of a particle with reference position $\tenq{x}\in R$. \\
Then the spatial (Eulerian) velocity or acceleration as 
\[
\left.
\begin{aligned}
 \tenq{q}(\tenq{y},t) = &\bar{\tenq{v}}(\tenq{y},t) ={\left.{ \dot{\tenq{y}}(\tenq{x},t)}\right|_{\tenq{x}=\hat{\tenq{x}}(\tenq{y},t)}}\\
 &\bar{\tenq{a}}(\tenq{y},t) ={\left.{ \ddot{\tenq{y}}(\tenq{x},t)}\right|_{\tenq{x}=\hat{\tenq{x}}(\tenq{y},t)}}
\end{aligned}\right\}\Rightarrow
\left.
\begin{aligned}
& \bar{\tenq{v}}(\hat{\tenq{y}}(\tenq{x},t),t) = \dot{\tenq{y}}(\tenq{x},t)\\
& \bar{\tenq{a}}(\hat{\tenq{y}}(\tenq{x},t),t) = \ddot{\tenq{y}}(\tenq{x},t)
\end{aligned}\right.
\]

\paragraph{Caution}\mbox{}\\

In general 
\[\frac{\partial}{\partial t}\bar{\tenq{v}}(\tenq{y},t)\neq\tenq{a}(\tenq{y},t)\]
However
\[{\left.{\frac{\partial}{\partial t}\left\{\bar{\tenq{v}}(\hat{\tenq{y}}(\tenq{x},t),t)\right\}}\right|_{\tenq{x}=\hat{\tenq{x}}(\tenq{y},t)}} = \bar{\tenq{a}}(\tenq{y},t)\]
This motivates the notion of material time derivative. 

\section{Material time derivative}

Consider a motion $\hat{\tenq{y}}\colon R \to R_t, t\in\tau$. Let $\psi(\tenq{y},t)$ be a spatial field and assume $\Psi(\tenq{x},t)$ is the corresponding referential (material) field, i.e. $\Psi(\tenq{x},t) = \psi(\hat{\tenq{y}}(\tenq{x},t),t)$. Then the material time derivative of $\psi$ is an other spatial field $\dot{\psi}(\tenq{y},t)$ such that
\begin{align*}
\dot{\psi}\footnotemark(\tenq{y},t) 
&= \left.{\left[\frac{\partial}{\partial t}\Psi(\tenq{x},t)\right]}\right|_{\tenq{x}=\hat{\tenq{x}}(\tenq{y},t)}\\
&= \left.{\left[\frac{\partial}{\partial t}\psi(\hat{\tenq{y}}(\tenq{x},t),t)\right]}\right|_{\tenq{x}=\hat{\tenq{x}}(\tenq{y},t)}
\end{align*}
\footnotetext{familiar notation $\frac{D\psi}{Dt}$}
\begin{notation}\mbox{}\\
\begin{enumerate}[label=(\arabic*)]
\item 
$\frac{\partial}{\partial t}\psi(\tenq{y},t)\neq\dot{\psi}(\tenq{y},t)$
\item 
$\dot{\Psi}(\tenq{x},t) = \frac{\partial}{\partial t}\Psi(\tenq{x},t)$
\end{enumerate}
\end{notation}

\begin{lemma}\mbox{}\\ \label{material derivative lemma-1}

Let $\phi(\tenq{y},t)$ and $\tenq{g}(\tenq{y},t)$ be two spatial fields (scalar, vector). \\
Then 

\[\frac{D\phi}{Dt} = \dot{\phi}(\tenq{y},t) = \frac{\partial}{\partial t}\phi(\tenq{y},t) + \nabla\phi(\tenq{y},t)\cdot\bar{\tenq{v}}(\tenq{y},t)\]
Similarly
\[\frac{D\tenq{g}}{Dt} = \dot{\tenq{g}}(\tenq{y},t) = \frac{\partial}{\partial t}\tenq{g}(\tenq{y},t) + \nabla\tenq{g}(\tenq{y},t)\bar{\tenq{v}}(\tenq{y},t)\]

\end{lemma}

\begin{proof}\mbox{}\\

Let $\Phi(\tenq{x},t) = \phi(\hat{\tenq{y}}(\tenq{x},t),t)$\\
Then 
\begin{align*}
\frac{\partial}{\partial t}\Phi(\tenq{x},t)
&= \left[\frac{\partial}{\partial t}\phi(\hat{\tenq{y}}(\tenq{x},t),\tau)\right]_{\tau=t} + \left[\frac{\partial}{\partial t}\phi(\hat{\tenq{y}}(\tenq{x},\tau),t)\right]_{\tau=t}\\
&= \left[\frac{\partial}{\partial t}\phi(\hat{\tenq{y}}(\tenq{x},\tau),t)\right]_{\tau=t} + \left.{\nabla \phi(\tenq{y},t)}\right|_{\tenq{y} = \hat{\tenq{y}}(\tenq{x},t)}\cdot\frac{\partial\hat{\tenq{y}}}{\partial t}
\end{align*}

\[\frac{\partial}{\partial t}\Phi(\tenq{x},t) = \frac{\partial}{\partial t}\phi(\tenq{y},t) + \nabla\phi(\tenq{y},t)\cdot\bar{\tenq{v}}(\tenq{y},t)\]

Let $\tenq{x}=\hat{\tenq{x}}(\tenq{y},t)$

\[\dot{\phi}(\tenq{y},t) = \frac{\partial}{\partial t}\phi(\tenq{y},t) + \nabla\phi(\tenq{y},t)\cdot\bar{\tenq{v}}(\tenq{y},t)\]


\begin{align*}
\frac{\partial}{\partial t}\left[g_i(\hat{\tenq{y}}(\tenq{x},t),t)\right]
&= \frac{\partial g_i}{\partial y_j}\frac{\partial y_j(\tenq{x},t)}{\partial t} + \frac{\partial g_i}{\partial t}\\
&= \left(\nabla \tenq{g}\right)_{ij}v_j + \frac{\partial g_i}{\partial t}
\end{align*}

\[\therefore \dot{\tenq{g}}(\hat{\tenq{y}}(\tenq{x},t),t) = \left(\nabla\tenq{g}\right)\bar{\tenq{v}} + \frac{\partial\tenq{g}(\hat{\tenq{y}}(\tenq{x},t),t)}{\partial t}\]

\end{proof}

\begin{lemma}\mbox{}\\

Let $\tenq{\tau}\in C^1[t_0,t_1]$ be a function whose values are \emph{non-singular} second order tensors. Then


\begin{equation}\tag{$\bigstar$} \label{material derivative -star}
\begin{split}
\frac{d}{dt}\det\tenq{\tau}(t)
&= \det\tenq{\tau}(t)\left\{\dot{\tenq{\tau}}(t)\cdot\tenq{\tau}^{-T}(t)\right\}\\
&= \det\tenq{\tau}(t)\dot{\tau}_{ij}\tau_{ji}^{-1}(t)\phantom{-}(t_0\leq t\leq t_1)
\end{split}
\end{equation}


\end{lemma}

\begin{proof}\mbox{}\\

Recall 

\[\det\tenq{\tau}(t) = \frac{1}{6}\epsilon_{ijk}\epsilon_{pqr}\tau_{ip}(t)\tau_{jq}(t)\tau_{kr}(t)\]

\begin{equation}\tag{1}\label{material derivative-1}
\begin{split}
\therefore \frac{d}{dt}\left\{\det\tenq{\tau}(t)\right\} 
&= \frac{1}{6}\epsilon_{ijk}\epsilon_{pqr}\left\{\dot{\tau}_{ip}(t)\tau_{jq}(t)\tau_{kr}(t)\right. \\
 &+ \left.\tau_{ip}(t)\dot{\tau}_{jq}(t)\tau_{kr}(t) + \tau_{ip}(t)\tau_{jq}(t)\dot{\tau}_{kr}(t)\right\}
\end{split}
\end{equation}



Recall also that

\begin{equation}\tag{2}\label{material derivative-2}
\boxed{\left(\tenq{\tau}^{-1}\right)_{ij}\equiv\tau_{ij}^{-1} = \frac{1}{2det\tenq{\tau}}\epsilon_{ipq}\epsilon_{jrs}\tau_{rp}\tau_{sq}}
\end{equation}

From \eqref{material derivative-1} and \eqref{material derivative-2}, we obtain 

\begin{align*}
\frac{d}{dt}\left\{\det\tenq{\tau}(t)\right\} 
&= \frac{1}{6}
\left\{
2\det\tenq{\tau}\dot{\tau}_{ip}\tau_{pi}^{-1} +
2\det\tenq{\tau}\dot{\tau}_{jq}\tau_{qj}^{-1}
\right. \\
&+ \left. 2\det\tenq{\tau}\dot{\tau}_{kr}\tau_{rk}^{-1}
\right\}
\end{align*}
or 

\[\frac{d}{dt}\left\{\det\tenq{\tau}(t)\right\} = \det\tenq{\tau}(t)\dot{\tau}_{ij}(t)\tau_{ji}^{-1}(t)\]


\end{proof}



\begin{claim}\mbox{}\\

Let $\hat{\tenq{y}}(\bullet,t)\colon R\to R_t$, one has the identity

\begin{equation}
\boxed{\dot{J}(\tenq{x},t) = J(\tenq{x},t)\nabla_{\tenq{y}}\cdot \tenq{q}(\tenq{y},t), \bar{\tenq{v}} = \tenq{q}(\tenq{y},t)} \tag{$\mathsection$} \label{material derivative-mathsection}
\end{equation}

\end{claim}

\begin{proof}\mbox{}\\

Let $\tenq{T}(t) = \nabla \hat{\tenq{y}}(\tenq{x},t)$, $\det \tenq{T} = \det\nabla \hat{\tenq{y}}(\tenq{x},t) = J(\tenq{x},t)$

\[\frac{d}{dt}\left\{\det\nabla \hat{\tenq{y}}(\tenq{x},t)\right\} = \dot{J}(\tenq{x},t)\vc{=}{\eqref{material derivative -star}}{\footnotemark} J(\tenq{x},t)\dot{\hat{y}}_{i,j}(\tenq{x},t)\hat{x}_{j,i}(\tenq{y},t)\]

\footnotetext{$\frac{d}{dt},\frac{\partial}{\partial x_i}$ interchangeable}

recall $\nabla\hat{\tenq{x}}(\tenq{y},t) = \left[\nabla \hat{\tenq{y}}(\tenq{x},t)\right]^{-1}$

\begin{align*}
\dot{J}(\tenq{x},t)
&\vc{=}{\footnotemark}{} J(\tenq{x},t)\frac{\partial}{\partial x_j}v_i(\tenq{x},t)\frac{\partial \hat{x}_j(\tenq{y},t)}{\partial y_i}\\
&\vc{=}{\footnotemark}{}J(\tenq{x},t)\frac{\partial}{\partial y_i}\bar{v}_i(\tenq{y},t)\\
&= J(\tenq{x},t)q_{i,i}\\
&= J(\tenq{x},t)\nabla_{\tenq{y}}\cdot\tenq{q}(\tenq{y},t)
\end{align*}

\footnotetext{$\frac{\partial \hat{\tenq{y}}}{\partial t} = \tenq{v}$ the Lagrangian velocity}
\footnotetext{Chain rule and $\tenq{q}(\tenq{y},t) = \bar{\tenq{v}}(\tenq{y},t) = \tenq{v}(\tenq{x}(\tenq{y},t),t)$}


\end{proof}

\begin{claim}\mbox{}\\

For any admissible motion mass balance (local) implies

\[\frac{\partial}{\partial t}\rho(\tenq{y},t) + \nabla\rho(\tenq{y},t)\cdot\bar{\tenq{v}}(\tenq{y},t) + \rho(\tenq{y},t)\nabla\cdot\bar{\tenq{v}}(\tenq{y},t) = 0\]

or 

\[\frac{\partial}{\partial t}\rho(\tenq{y},t) + \nabla \cdot\left[\rho(\tenq{y},t)\bar{\tenq{v}}(\tenq{y},t)\right] = 0\]



\end{claim}

\begin{proof}\mbox{}\\

Recall $\eqref{mass conservation-**}$ and $\eqref{material derivative-mathsection}$
\[\rho(\tenq{y},t)J(\tenq{x},t)=\rho_0(\tenq{x})\]
\[\dot{J}(\tenq{x},t) = J(\tenq{x},t)\nabla_{\tenq{y}}\cdot \bar{\tenq{v}}(\tenq{y},t)\]
Differentiate with respect to $t$ ($x$ fixed)

\begin{align*}
&\rightarrow
\left[\frac{\partial}{\partial t}\rho(\hat{\tenq{y}}(\tenq{x},t),t)\right]J(\tenq{x},t) + \rho(\hat{\tenq{y}}(\tenq{x},t),t)\frac{\partial}{\partial t}J(\tenq{x},t) = 0\\
&\rightarrow
\dot{\rho}(\tenq{y},t)J(\tenq{x},t) + \rho(\tenq{y},t)J(\tenq{x},t)\nabla\cdot\bar{\tenq{v}}(\tenq{y},t) = 0\\
&\rightarrow
\dot{\rho}(\tenq{y},t) + \rho(\tenq{y},t)\nabla\cdot\bar{\tenq{v}}(\tenq{y},t) = 0\\
&\vc{\rightarrow}{Lemma}{\ref{material derivative lemma-1}}
\frac{\partial}{\partial t}\rho(\tenq{y},t) + \nabla\rho(\tenq{y},t)\cdot\bar{\tenq{v}}(\tenq{y},t) + \rho(\tenq{y},t)\nabla\cdot\bar{\tenq{v}}(\tenq{y},t) = 0\\
&\rightarrow
\frac{\partial}{\partial t}\rho(\tenq{y},t) + \nabla \cdot\left[\rho(\tenq{y},t)\bar{\tenq{v}}(\tenq{y},t)\right] = 0
\end{align*}



\end{proof}

\section{For locally volume preserving (isochoric motion)}

\[J(\tenq{x},t)=1\]
\begin{align*}
&\Rightarrow \dot{J}(\tenq{x},t)=0 \vc{=}{\eqref{material derivative-mathsection}}{}J(\tenq{x},t)\nabla\cdot\bar{\tenq{v}}(\tenq{y},t) = 0\\
&\Rightarrow\nabla\cdot\bar{\tenq{v}} = 0\\
&\Rightarrow \dot{\rho} (\tenq{y},t) = 0
\end{align*}
or 
\[\frac{\partial}{\partial t}\rho(\tenq{y},t) + \nabla\rho(\tenq{y},t)\cdot\bar{\tenq{v}}(\tenq{y},t)= 0\]

\section{Spatial velocity gradient tensor}

Let $\bar{\tenq{v}}(\tenq{y},t)$ be the spatial or Eulerian velocity field. Define a tensor field $\tenq{L}(\bullet,t)\colon R_t \in \mathscr{L}$\footnote{$\mathscr{L}$ is set of all $2^{nd}$ order tensor} by 

\[\tenq{L}(\tenq{y},t) = \nabla_{\tenq{y}}\bar{\tenq{v}}(\tenq{y},t)\phantom{-}\forall\tenq{y}\in R_t\phantom{-}\forall t\in \tau\]
called the spatial velocity gradient tensor. 

Recall 

\begin{align*}
\tenq{F}(\tenq{x},t) 
&= \nabla_{\tenq{x}}\hat{\tenq{y}}(\tenq{x},t)\\
\Rightarrow \dot{\tenq{F}}(\tenq{x},t) 
&= \vc{\left[\nabla_{\tenq{x}}\hat{\tenq{y}}(\tenq{x},t)\right]}{\bullet}{} = \nabla_{\tenq{x}}\dot{\hat{\tenq{y}}}(\tenq{x},t)
\end{align*}
However recall 

\begin{align*}
\dot{\hat{\tenq{y}}}(\tenq{x},t) 
&= \bar{\tenq{v}}(\hat{\tenq{y}}(\tenq{x},t),t)\\
\Rightarrow \nabla_{\tenq{x}}\dot{\hat{\tenq{y}}}(\tenq{x},t)
&= \left.\nabla_{\tenq{y}}\bar{\tenq{v}}(\tenq{y},t)\right|_{\tenq{y} = \hat{\tenq{y}}(\tenq{x},t)}\nabla_{\tenq{x}}\hat{\tenq{y}}(\tenq{x},t)
\end{align*}

\[\therefore\boxed{\dot{\tenq{F}}(\tenq{x},t) = \tenq{L}(\hat{\tenq{y}}(\tenq{x},t),t)\tenq{F}(\tenq{x},t)}\]

\begin{definition}\mbox{}\\

The stretching and spin tensor fields as 

\begin{align*}
\tenq{D}(\tenq{y},t)
&= \textit{sym } \tenq{L}(\tenq{y},t) = \frac{1}{2}\left[\nabla \bar{\tenq{v}}(\tenq{y},t) + \left(\nabla\bar{\tenq{v}}(\tenq{y},t)\right)^T\right]\\
\tenq{W}(\tenq{y},t)
&= \textit{skew } \tenq{L}(\tenq{y},t) = \frac{1}{2}\left[\nabla \bar{\tenq{v}}(\tenq{y},t) - \left(\nabla\bar{\tenq{v}}(\tenq{y},t)\right)^T\right]
\end{align*}

so that $\tenq{L} = \tenq{W} + \tenq{D}$, $\tenq{D} = \tenq{D}^T$, $\tenq{W} = -\tenq{W}^T$

\end{definition}


\begin{note}\mbox{}\\

Recall $\tenq{F} = \tenq{Q}\tenq{U} = \tenq{V}\tenq{U}$ ($F$, $Q$, $U$, and $V$ are all functions of time) and for infinitesimal deformation we have

\begin{align*}
\tenq{L} &\sim \dot{\tenq{F}}\\
\tenq{D} &\sim \dot{\tenq{U}}\\
\tenq{W} &\sim \dot{\tenq{Q}}\\
\end{align*}

\end{note}

Consider two particles $P$ and $Q$ with material position $\tenq{x}$ and $\tenq{x}+d\tenq{x}$. At time $t$ they are to be found at $\tenq{y}(\tenq{x},t)$ and $\tenq{y}(\tenq{x}+d\tenq{x},t)$

\begin{figure}[H]
\centering
\resizebox{0.7\textwidth}{!}{\input{figure25.eps_tex}}
\caption{Scheme of deformation of $P$ and $Q$}
\label{Scheme of deformation of P and Q}
\end{figure}

\[\tenq{y}(\tenq{x}+d\tenq{x},t) = \tenq{y}(\tenq{x},t) + \nabla_{\tenq{x}}\tenq{y}(\tenq{x},t)d\tenq{x} + \text{ higer order term (H.O.T.)}\]

\[\therefore d\tenq{y} = \underbrace{\nabla_{\tenq{x}}\tenq{y}}_{\tenq{F}}d\tenq{x} \text{ or }dy_i = \frac{\partial y_i}{\partial x_j}dx_j\]

Velocity of $P$\\

\[\tenq{v} = \dot{\tenq{y}}(\tenq{x},t)\]
Relative velocity of $P'$ and $Q'$

\begin{align*}
d\tenq{v} = \vc{\left(d\tenq{y}\right)}{\bullet}{} 
&= \vc{\left(\nabla_{\tenq{x}}\tenq{y}\right)}{\bullet}{}d\tenq{x}\\
&= \nabla_{\tenq{x}}\tenq{v}(\tenq{x},t)d\tenq{x}\\
&= \frac{\partial v_i(\tenq{x},t)}{\partial x_j}dx_j\tenq{e}_i\\
&= \frac{\partial \bar{v}_i(\hat{\tenq{y}}(\tenq{x},t),t)}{\partial y_m}\underbrace{\frac{\partial y_m(\tenq{x},t)}{\partial x_j}dx_j}_{dy_m}\tenq{e}_i\\
&= \underbrace{\left(\nabla_{\tenq{y}}\bar{\tenq{v}}\right)}_{\text{velocity}\atop \text{gradient tensor}}d\tenq{y}
\end{align*}

where relative velocity in terms of \emph{current} relative position. 
\begin{equation}
\Rightarrow dv_i = \frac{\partial v_i}{\partial x_j}dx_j = \frac{\partial \bar{v}_i}{\partial y_j}dy_j \tag{$\blacktriangle$} \label{Spatial velocity gradient tensor-triangle}
\end{equation}


The stretching (or rate of strain) and spin tensors are defined as 

\begin{align*}
\tenq{D}
&\vc{=}{\Delta}{}\textit{sym } \nabla_{\tenq{y}}\bar{\tenq{v}} = \frac{1}{2}\left[\nabla\bar{\tenq{v}} + \left(\nabla\bar{\tenq{v}}\right)^T\right]\\
\tenq{W}
&\vc{=}{\Delta}{}\textit{skew } \nabla_{\tenq{y}}\bar{\tenq{v}} = \frac{1}{2}\left[\nabla\bar{\tenq{v}} - \left(\nabla\bar{\tenq{v}}\right)^T\right]
\end{align*}

\begin{note}

Recall that exists a vector $\tenq{\omega}$ such that

\[\omega_i = -\frac{1}{2}\epsilon_{ijk}W_{jk}\]

\end{note}


\section{Physical interpretation of \texorpdfstring{$D_{ij}$}{Lg}}

Let the length of the line segment $P'Q'$ is $ds$, then 

\[ds^2 = dy_idy_i = \frac{\partial y_i}{\partial x_j}\frac{\partial y_i}{\partial x_k}dx_jdx_k\]

\begin{align*}
\vc{\left(ds^2\right)}{\bullet}{}
&= \left(\frac{\partial \dot{y}_i}{\partial x_j}\frac{\partial y_i}{\partial x_k} + \frac{\partial y_i}{\partial x_j}\frac{\partial \dot{y}_i}{\partial x_k}\right)dx_jdx_k\\
&= \left(\frac{\partial v_i}{\partial x_j}\frac{\partial y_i}{\partial x_k} + \frac{\partial v_i}{\partial x_k}\frac{\partial y_i}{\partial x_j}\right)dx_jdx_k\\
&\vc{=}{\eqref{Spatial velocity gradient tensor-triangle}}{}\frac{\partial \bar{v}_i}{\partial y_j}dy_jdy_i + \frac{\partial \bar{v}_i}{\partial y_k}dy_kdy_i\\
&= 2\frac{\partial \bar{v}_i}{\partial y_j}dy_idy_j\\
&= 2\left(D_{ij} + W_{ij}\right)dy_idy_j = \footnotemark 2D_{ij}dy_idy_j
\end{align*}

\footnotetext{$W_{ij} = -W_{ji}$}

\[\therefore 2ds\vc{\left(ds\right)}{\bullet}{} = 2D_{ij}dy_idy_j\]

or 

\begin{equation}
\frac{1}{ds}\vc{\left(ds\right)}{\bullet}{}  = D_{ij}\frac{dy_i}{ds}\frac{dy_j}{ds} \cdot \text{rate of stretching} \tag{$\ast$} \label{Physical interpretation of D_ij-*}
\end{equation}

Let $\frac{d\tenq{y}}{ds} = \tenq{e}_1$, then $D_{11}$ represents the rate of strain of a fluid element parallel to $\tenq{e}_1$ axis. Similar to $D_{22}$ and $D_{33}$. 

\begin{note}\mbox{}\\

\[D_{11} + D_{22} + D_{33} = \text{rate of volume dilatation}\]

Again consider two segment $PQ$ and $PR$, $R$ is the particle with material position $\tenq{x}+d\tenq{x}'$

\begin{figure}[H]
\centering
\resizebox{0.7\textwidth}{!}{\input{figure26.eps_tex}}
\caption{Scheme of change of two segment}
\label{Scheme of change of two segment}
\end{figure}

\[dsds'\cos\theta = dy_idy_i'\]

Take material time derivative 

\begin{align*}
\vc{\left(dsds'\cos\theta\right)}{\bullet}{} 
&= dv_idy_i' + dv_i'dy_i\\
&\vc{=}{\eqref{Spatial velocity gradient tensor-triangle}}{} \frac{\partial \bar{v}_i}{\partial y_j}dy_jdy_i' + \underbrace{\frac{\partial \bar{v}_i}{\partial y_j}dy_j'dy_i}_{i\leftrightarrow j}
\end{align*}

\begin{align*}
&\Rightarrow
\cos\theta \left[\vc{\left(ds\right)}{\bullet}{}ds' + ds\vc{\left(ds'\right)}{\bullet}{}\right] - \vc{\left(\theta\right)}{\bullet}{}\sin\theta dsds' = \left(\frac{\partial \bar{v}_i}{\partial y_j} + \frac{\partial \bar{v}_j}{\partial y_i}\right)dy_i'dy_j\\
&\Rightarrow
\cos\theta\left[\frac{1}{ds}\vc{\left(ds\right)}{\bullet}{} + \frac{1}{ds'}\vc{\left(ds'\right)}{\bullet}{}\right]-\dot{\theta}\sin\theta = 2D_{ij}\frac{dy_i'}{ds'}\frac{dy_j}{ds}
\end{align*}

Let $\frac{d\tenq{y}'}{ds'} = \tenq{e}_1$, $\frac{d\tenq{y}}{ds} = \tenq{e}_2$\\
then $\theta = \frac{\pi}{2}$

\[\therefore -\dot{\theta} = 2D_{12}\]

$D_{12}$ is one-half the rate of decrease of the angle between two segments originally parallel to $\tenq{e}_1$ and $\tenq{e}_2$ respectively. 

\end{note}


\section{Conservation of mass for a control volume}

If $D$ is a control volume independent of time, then 

\[\frac{d}{dt}\int_D \rho(\tenq{y},t)d\volume_{\tenq{y}} + \int_{\partial D} \rho(\tenq{y},t) \bar{\tenq{v}}(\tenq{y},t)\cdot\tenq{n}(\tenq{y},t)dA_{\tenq{y}} = 0\]

\begin{figure}[H]
\centering
\resizebox{0.5\textwidth}{!}{\input{figure27.eps_tex}}
\caption{Scheme of control volume}
\label{Scheme of control volume}
\end{figure}

\begin{proof}\mbox{}\\

Recall $D$ is time independent 

\[\int_D \frac{\partial}{\partial t}\rho(\tenq{y},t)d\volume_{\tenq{y}} = \frac{d}{dt}\int_D \rho(\tenq{y},t)d\volume_{\tenq{y}}\]

also 

\[\int_D \nabla\cdot\left[\rho\bar{\tenq{v}}\right]d\volume_{\tenq{y}} = \int_{\partial D}\rho\bar{\tenq{v}}\cdot\tenq{n}dA_{\tenq{y}}\phantom{-}\text{(Divergence theorem)}\]

so 

\[\int_D\left\{\frac{\partial}{\partial t}\rho(\tenq{y},t) + \nabla \cdot \left[\rho(\tenq{y},t)\bar{\tenq{v}}(\tenq{y},t)\right]\right\}d\volume_{\tenq{y}} = 0\]

\end{proof}

\section{Reynolds transport theorem (RTT)}
Let $P\in R$, $P_t \in R_t$, $P_t = \hat{\tenq{y}}(P,t)$\\
and define 

\[\phi(t) = \int_{P_t}\rho(\tenq{y},t)\psi(\tenq{y},t)d\volume_{\tenq{y}}\]
Then 

\[\frac{d}{dt}\phi(t) = \int_{P_t}\rho(\tenq{y},t)\dot{\psi}(\tenq{y},t)d\volume_{\tenq{y}}\]
or 

\[\frac{d\phi}{dt} = \int_{P_t}\rho\left[\frac{\partial\psi}{\partial t} + \nabla \psi\cdot\bar{\tenq{v}}\right]d\volume_{\tenq{y}}\]

\begin{proof}\mbox{}\\

\begin{align*}
\phi(t) 
&= \int_P \rho(\hat{\tenq{y}}(\tenq{x},t),t)\psi(\hat{\tenq{y}}(\tenq{x},t),t)J(\tenq{x},t)d\volume_{\tenq{x}}\\
&\vc{=}{\eqref{mass conservation-**}}{}\int_P\rho_0(\tenq{x})\psi(\hat{\tenq{y}}(\tenq{x},t),t)d\volume_{\tenq{x}}\\
\end{align*}

\begin{align*}
\frac{d\phi}{dt} 
&= \int_P\rho_0(\tenq{x})\frac{\partial}{\partial t}\psi(\hat{\tenq{y}}(\tenq{x},t),t)d\volume_{\tenq{x}}\\
&= \int_P\left.\left[\rho(\tenq{y},t)\dot{\psi}(\tenq{y},t)\right]\right|_{\tenq{y} = \hat{\tenq{y}}(\tenq{x},t)}J(\tenq{x},t)d\volume_{\tenq{x}}\\
\end{align*}

\[\therefore \frac{d\phi}{dt} = \int_{P_t}\rho(\tenq{y},t)\dot{\psi}(\tenq{y},t)d\volume_{\tenq{y}}\]

\end{proof}

If $\Psi(\tenq{x},t) = \psi(\hat{\tenq{y}}(\tenq{x},t),t)$ then 

\[\frac{d\phi}{dt} = \int_P\rho_0(\tenq{x})\frac{\partial}{\partial t}\Psi(\tenq{x},t)d\volume_{\tenq{x}}\]

\section{Body forces, contact forces, surface tractions}

Let $\hat{\tenq{y}}(\bullet,t)\colon R \to R_t$. Let $P\subset R, P_t = \hat{\tenq{y}}(P,t)$. \\

Admissible forces:

\begin{enumerate}[label=(\alph*)]

\item
body forces (extrinsic, assigned), e.g. gravity, magnetic field. 

\item contact forces. (contact of material in $P_t$ with surrounding material)

\end{enumerate}

\begin{figure}[H]
\centering
\resizebox{0.4\textwidth}{!}{\input{figure28.eps_tex}}
\caption{Scheme of deformation of $P$}
\label{Scheme of deformation of P_1}
\end{figure}

\subsection{Body force density}

Assume that exists a vector valued function $\tenq{b}(\bullet, t)$ defined on $R_t$ called the body force per unit current volume such that

\begin{enumerate}[label=(\roman*)]

\item
$\tenq{b}(\bullet,t)\in C(R_t)$
\item
The resultant force and resultant moment about $O$ due to the body forces acting on the material occupying $R_t$ are given by 

\[\int_{R_t}\tenq{b}(\tenq{y},t)d\volume_{\tenq{y}},\int_{R_t}\tenq{y}\wedge\tenq{b}(\tenq{y},t)d\volume_{\tenq{y}}\]
dim $\tenq{b} = [FL^{-3}]$
\end{enumerate}


\subsection{Contact forces, surface tractions}

Assume there is a vector valued function $\tenq{t}(\bullet,\bullet,t)$ defined on $R_t\times u$ ($t_0\leq t \leq t_1$), $u\cdots$ set of all unit vector, called the surface force per unit current area or the actual surface traction or Cauchy traction with the properties: 

\begin{enumerate}[label=(\roman*)]

\item
$\tenq{t}(\bullet,\bullet,t)\in C(R_t\times u)$
\item
If $S$ is a regular subsurface of $P$ and if $S_t = \hat{\tenq{y}}(S,t)\phantom{-}\forall t\in\tau$, then the resultant force and resultant moment about $O$ due to the contact forces acting over $S_t$ on the material occupying $P_t$ at time t are given by 

\[\int_{S_t}\tenq{t}(\tenq{y},\tenq{n}(\tenq{y}),t)dA_{\tenq{y}}\cdots\text{resultant force}\]
\[\int_{S_t}\tenq{y}\wedge\tenq{t}(\tenq{y},\tenq{n}(\tenq{y}),t)dA_{\tenq{y}}\cdots\text{resultant moment}\]
dim $\tenq{t} = [FL^{-2}]$

\begin{figure}[H]
\centering
\resizebox{0.42\textwidth}{!}{\input{figure29.eps_tex}}
\caption{Scheme of contact forces, surface tractions}
\label{Scheme of contact forces, surface tractions}
\end{figure}

\end{enumerate}



\section{Conservation of linear momentum}

\begin{align*}
\int_{\partial P_t}\tenq{t}(\tenq{y},\tenq{n}(\tenq{y}),t)dA_{\tenq{y}} + \int_{P_t}\tenq{b}(\tenq{y},t)d\volume_{\tenq{y}}
&= \frac{d}{dt} \int_{P_t}\rho(\tenq{y},t)\bar{\tenq{v}}(\tenq{y},t)d\volume_{\tenq{y}}\\
&\vc{=}{RTT}{}\int_{P_t}\rho(\tenq{y},t)\bar{\tenq{a}}(\tenq{y},t)d\volume_{\tenq{y}}
\end{align*}

\section{Conservation of angular momentum}


\begin{align*}
&\int_{\partial P_t}\tenq{y}\wedge \tenq{t}(\tenq{y},\tenq{n}(\tenq{y}),t)dA_{\tenq{y}} 
+ \int_{P_t}\tenq{y}\wedge \tenq{b}(\tenq{y},t)d\volume_{\tenq{y}}\\
&= \frac{d}{dt} \int_{P_t}\rho(\tenq{y},t)\tenq{y}\wedge \bar{\tenq{v}}(\tenq{y},t)d\volume_{\tenq{y}}\\
&\vc{=}{RTT}{}\int_{P_t}\rho\vc{[\tenq{y}\wedge\bar{\tenq{v}}]}{\bullet}{} d\volume_{\tenq{y}}\\
&\vc{=}{ \footnotemark}{}\int_{P_t}\rho\tenq{y}\wedge\bar{\tenq{a}}(\tenq{y},t)d\volume_{\tenq{y}}
\end{align*}

\footnotetext{$\dot{\tenq{y}}\wedge\bar{\tenq{v}} = \bar{\tenq{v}}\wedge\bar{\tenq{v}} = \tenq{0}$}

\begin{example}{For $\tenq{b}(\tenq{y},t)$}

Let the center of the earth be $O$, $O\not\in R_t$, mass of earth $M$, then the gravitational body force density 

\[\tenq{b}(\tenq{y},t) = -\frac{\rho(\tenq{y},t)GM\tenq{y}}{\abs{\tenq{y}}^3}\phantom{-}\forall\tenq{y}\in R_t, t\in \tau\]

where $G>0\cdots$ Universal constant of Gravitation.

\end{example}

\section{Existence of stress}

The following celebrated theorem by Cauchy introduces and establishes the notion of stress based on Cauchy's postulate that \emph{tractions depend on the unit normal}, and by using linear momentum balance. 

\begin{theorem}{Cauchy}

Let $\hat{\tenq{y}}(\bullet,t)\colon R \to R_t$ be a motion and let $\tenq{t}(\bullet,\bullet,t)$ and $\tenq{b}(\bullet,t)$ be a traction and a body force density field such that linear momentum balance holds. Then there exists a second order tensor field $\tenq{\tau}(\bullet,t)\colon R_t \in \mathscr{L}$\footnote{$\mathscr{L}$ is set of all $2^{nd}$ order tensor} for each $t\in[t_0,t_1]$ such that

\[\boxed{\tenq{t}(\tenq{n},\tenq{y},t) = \tenq{\tau}(\tenq{y},t)\tenq{n}}\phantom{-}\forall \tenq{n}\in u,\forall \tenq{y}\in R_t,\forall t\in \tau\]

or 

\[t_i^X(\tenq{n},\tenq{y},t) = \tau_{ij}^X(\tenq{y},t)n_j^X\]

\end{theorem}

\begin{proof}\mbox{}\\

Fix $t\in [t_0,t_1]$

\paragraph{Assertion 1}\mbox{}\\

Given any $\tenq{y}\in R_t$ for $X = \left\{ {0;\tenq{e}_1,\tenq{e}_2,\tenq{e}_3} \right\} \in \mathscr{F}$ and unit vector $\tenq{n}$ with $\tenq{n}\cdot\tenq{e}_i>0$\\
then 

\[\tenq{t}(\tenq{n},\tenq{y}) = -\sum_{i=1}^{3}n_i\tenq{t}(-\tenq{e}_i,\tenq{y}), \phantom{-}n_i = \tenq{n}\cdot\tenq{e}_i\]
\begin{proof}\mbox{}\\

\begin{figure}[H]
\centering
\resizebox{1.0\textwidth}{!}{\input{figure30.eps_tex}}
\caption{Scheme of tetrahedra}
\label{Scheme of tetrahedra}
\end{figure}

\emph{Fix} $\tenq{y} = \mathring{\tenq{y}}$ in $R_t$ ($\mathring{\tenq{y}}$ is arbitrary)\\

Let $h>0$ and consider the $h$-parameter family of tetrahedra $P(h)$
\[
\pi_i(h) \text{ has } -\tenq{e}_i \text{ as normal }
\]

\[
\pi(h) \text{ has } \tenq{n} \text{ as normal }
\]

Area $\pi(h) = A(h) = c_1h^2$, where $c_1$ is a real number\\
Area $\pi_i(h) = A_i(h) = A(h)n_i$

Linear momentum for $P(h)\Rightarrow$

\begin{equation}
\int_{\partial P(h)}\tenq{t}(\tenq{n},\tenq{y})dA_{\tenq{y}} = \int_{P(h)}\left[\rho\bar{\tenq{a}}-\tenq{b}\right]d\volume_{\tenq{y}}\phantom{-},\text{volume } P(h) = c_2h^3
\tag{1}\label{Existence of stress-(1)}
\end{equation}

the integral of R.H.S. of above is continuous on $R_t$ so exists $M$ such that $\abs{\rho\bar{\tenq{a}}-\tenq{b}}\leq M$ on $P(h)\phantom{-}\forall h\in [0,h_0]$

\[\Rightarrow \abs{\int_{P(h)}\left[\rho\bar{\tenq{a}}-\tenq{b}\right]d\volume_{\tenq{y}}}\leq M\underbrace{c_2h^3}_{\text{volume}}\]
from \eqref{Existence of stress-(1)} $\Rightarrow$

\[\abs{\frac{1}{A(h)}\int_{\partial P(h)}\tenq{t}(\tenq{n},\tenq{y})dA_{\tenq{y}}}\leq\frac{Mc_2h^3}{A(h)} = \frac{Mc_2h^3}{c_1h^2}\to 0\text{ as } h\to 0\]

\begin{equation}
\Rightarrow \frac{1}{A(h)}\int_{\partial P(h)}\tenq{t}(\tenq{n},\tenq{y})dA_{\tenq{y}}\to 0\text{ as } h\to 0
\tag{2}\label{Existence of stress-(2)}
\end{equation}
or 


\begin{multline*}
\frac{1}{A(h)}\int_{\partial P(h)}\tenq{t}(\tenq{n},\tenq{y})dA_{\tenq{y}} = \frac{1}{A(h)}\int_{\pi(h)}\tenq{t}(\tenq{n},\tenq{y})dA_{\tenq{y}} \\
+ \frac{1}{A(h)}\sum_{i=1}^3\int_{\pi_i(h)}\tenq{t}(-\tenq{e}_i,\tenq{y})dA_{\tenq{y}}
\end{multline*}

also since traction field is continuous we have that

\[\abs{\tenq{t}(\tenq{n},\tenq{y}) - \tenq{t}(\tenq{n},\mathring{\tenq{y}})}\leq m(h)\to 0 \text{ as } h\to 0\]

thus 

\[\frac{1}{A(h)}\abs{\int_{\pi(h)}\left[\tenq{t}(\tenq{n},\tenq{y}) - \tenq{t}(\tenq{n},\mathring{\tenq{y}})\right]dA}\leq\frac{m(h)A(h)}{A(h)}=m(h)\to 0\]
\[\text{ as } h\to 0\]

hence

\begin{align*}
\frac{1}{A(h)}\int_{\pi(h)}\tenq{t}(\tenq{n},\tenq{y})dA_{\tenq{y}}
&\to \frac{1}{A(h)}\tenq{t}(\tenq{n},\mathring{\tenq{y}})\underbrace{\int_{\pi(h)}dA_{\tenq{y}}}_{=A(h)}\text{ as } h\to 0\\
&\to \tenq{t}(\tenq{n},\mathring{\tenq{y}})\text{ as } h\to 0
\end{align*}

Similarly

\[\frac{A_i(h)}{A(h)}\frac{1}{A_i(h)}\int_{\pi_i(h)}\tenq{t}(-\tenq{e}_i,\tenq{y})dA_{\tenq{y}}\to \frac{A_i(h)}{A(h)}\tenq{t}(-\tenq{e}_i,\mathring{\tenq{y}})\]

where $A_i(h) = A(h)n_i$

\[\eqref{Existence of stress-(2)}\Rightarrow \tenq{t}(\tenq{n},\mathring{\tenq{y}}) + \sum_{i=1}^3n_i\tenq{t}(-\tenq{e}_i,\mathring{\tenq{y}}) = 0 \text{ for arbitrary }\mathring{\tenq{y}}\]

\begin{equation}
\therefore \tenq{t}(\tenq{n},\tenq{y}) = -\sum_{i=1}^3\tenq{t}(-\tenq{e}_i,\tenq{y})n_i\phantom{-}\forall \tenq{y}\in R_t
\tag{3}\label{Existence of stress-(3)}
\end{equation}





\end{proof}

Let $\tenq{n}\to \tenq{e}_1$, $\tenq{e}_2$, and $\tenq{e}_3$ respectively
\[
\left.
\begin{aligned}
& \tenq{t}(\tenq{e}_1,\tenq{y}) = -\tenq{t}(-\tenq{e}_1,\tenq{y})\\
& \tenq{t}(\tenq{e}_2,\tenq{y}) = -\tenq{t}(-\tenq{e}_2,\tenq{y})\\
& \tenq{t}(\tenq{e}_3,\tenq{y}) = -\tenq{t}(-\tenq{e}_3,\tenq{y})
\end{aligned}
\right\}
\tenq{t}(\tenq{n},\tenq{y}) = -\tenq{t}(-\tenq{n},\tenq{y})
\]
$\cdots$ Newton's law of action and reaction\\

So \eqref{Existence of stress-(3)} becomes

\begin{equation}
\tenq{t}(\tenq{n},\tenq{y}) = \sum_{i=1}^3n_i\tenq{t}(\tenq{e}_i,\tenq{y})\phantom{-}\forall\tenq{y}\in R_t\phantom{-}\forall\tenq{n}\in u
\tag{4}\label{Existence of stress-(4)}
\end{equation}

\fbox{
\begin{minipage}{25em}
If we repeat the above proof in other quadrants
, i.e. $n_i$ may not positive, the result still hold.
Thus on can remove the restriction $n_i>0$
\end{minipage}
}

Define the following 9 real numbers as follows

\begin{equation}
\tau_{ij}^X = \tenq{t}(\tenq{e}_j,\tenq{y})\cdot\tenq{e}_i
\tag{5}\label{Existence of stress-(5)}
\end{equation}

\fbox{
\begin{minipage}{25em}
This is not consistent with the convention 
we use in solid \\mechanics 
$\tau_{ij}^X = \tenq{t}(\tenq{e}_i,\tenq{y})\cdot\tenq{e}_j$
\end{minipage}
}
\begin{figure}[H]
\centering
\resizebox{0.4\textwidth}{!}{\input{figure31.eps_tex}}
\caption{Scheme of orientation of stress}
\label{Scheme of orientation of stress}
\end{figure}

e.g.
\[\tau_{11}^X = \tenq{t}(\tenq{e}_1,\tenq{y})\cdot\tenq{e}_1\]
\[\tau_{12}^X = \tenq{t}(\tenq{e}_2,\tenq{y})\cdot\tenq{e}_1\]

dot multiply \eqref{Existence of stress-(4)} by $\tenq{e}_i$ and use definition \eqref{Existence of stress-(5)}

\[t_i^X(\tenq{n},\tenq{y}) = \tenq{t}^X(\tenq{n},\tenq{y})\cdot\tenq{e}_i = \sum_{j=1}^3 \underbrace{\tenq{t}^X(\tenq{e}_j,\tenq{y})\cdot\tenq{e}_i}_{\tau_{ij}^X}n_j^X\]

\[\Rightarrow t_i^X(\tenq{n},\tenq{y}) = \tau_{ij}^Xn_j^X\]
or 
\[t_i^X(\tenq{n},\tenq{y}) = \tau_{ij}^X(\tenq{y})n_j^X\phantom{-}\forall X\in \mathscr{F},\forall\tenq{y}\in R_t, \forall \tenq{n}\in u\]

by quotient rule $\tau_{ij}^X$ represent component in $X$ of a tensor $\tenq{\tau}$

\[\Rightarrow \tenq{t}(\tenq{n},\tenq{y}) = \tenq{\tau}(\tenq{y})\tenq{n}\]

\[\tenq{\tau}(\tenq{y})\cdots\text{Cauchy stress tensor field}\]

\end{proof}

\begin{note}\mbox{}\\

In general, we release time so that 

\[\tenq{t}(\tenq{n}\footnotemark,\tenq{y},t) = \tenq{\tau}(\tenq{y},t)\tenq{n}\]

\end{note}

\footnotetext{$\tenq{n}(\tenq{y},t)$}

\section{Local version of balance laws}

\begin{theorem}\mbox{}\\

The traction field $\tenq{t}$ and body force density $\tenq{b}$ associated with the motion $\hat{\tenq{y}}(\bullet,t)\colon R\to R_t$ satisfy linear and angular momentum balance if 

\begin{enumerate}[label=(\roman*)]

\item
That exists $\tenq{\tau}(\bullet,t)\in C^1(R_t)$ such that $\tenq{t}(\tenq{n},\tenq{y},t) = \tenq{\tau}(\tenq{y},t)\tenq{n}$

\item
$\nabla \cdot \tenq{\tau}(\tenq{y},t) + \tenq{b}(\tenq{y},t) = \rho(\tenq{y},t)\bar{\tenq{a}}(\tenq{y},t)\phantom{-}\forall \tenq{y}\in R_t,\forall t\in \tau$

\item
$\tenq{\tau}(\tenq{y},t) = \tenq{\tau}^T(\tenq{y},t)$

\end{enumerate}

\end{theorem}

\begin{proof}\mbox{}\\

\[\int_{\partial P_t}\tenq{t}dA_{\tenq{y}} = \int_{\partial P_t}\tenq{\tau}\tenq{n}dA_{\tenq{y}} = \int_{P_t}\left[\rho\bar{\tenq{a}}-\tenq{b}\right]d\volume_{\tenq{y}}\]

by divergence theorem

\[\Rightarrow\int_{P_t}\left[\nabla\cdot\tenq{\tau}+\tenq{b}-\rho\bar{\tenq{a}}\right]d\volume_{\tenq{y}} = 0\phantom{-}\forall P_t\subset R_t\]

by continuity of integration and localization argument

\begin{equation}
\Rightarrow \nabla\cdot\tenq{\tau}(\tenq{y},t) + \tenq{b}(\tenq{y},t) = \rho(\tenq{y},t)\bar{\tenq{a}}(\tenq{y},t)\phantom{-}\forall\tenq{y}\in R_t
\tag{$\ast$}\label{Local version of balance laws-*}
\end{equation}

or 

\[\tau_{ij,j}(\tenq{y},t)+b_i(\tenq{y},t) = \rho(\tenq{y},t)\bar{a}_i(\tenq{y},t)\]

Also from conservation of angular momentum and linear momentum

\[\int_{\partial P_t}\tenq{y}\wedge\tenq{\tau}\tenq{n}dA_{\tenq{y}} + \int_{P_t}\tenq{y}\wedge\left[\tenq{b}-\rho\bar{\tenq{a}}\right]d\volume_{\tenq{y}} = \tenq{0}\]

Let $B_{il} = \epsilon_{ijk}y_j\tau_{kl}$,\footnotemark $\tenq{B} = B_{ij}\tenq{e}_i\otimes \tenq{e}_j$

\footnotetext{Define a vector $\tenq{g} = \tenq{\tau}\tenq{n}$ or $g_k = \tau_{kl}n_l$, then $\tenq{y}\wedge\tenq{\tau}\tenq{n} = \tenq{y}\wedge\tenq{g}$ or $\epsilon_{ijk}y_j\tau_{kl}n_l = \epsilon_{ijk}y_jg_k$. If we define $\tenq{B} = \tenq{y}\wedge\tenq{\tau}$ or $B_{il} = \epsilon_{ijk}y_j\tau_{kl}$ such that $\tenq{y}\wedge\tenq{\tau}\tenq{n} = \tenq{B}\tenq{n}$ or $\epsilon_{ijk}y_j\tau_{kl}n_l = B_{il}n_l$}

\[\Rightarrow \int_{\partial P_t}\tenq{B}\tenq{n}dA_{\tenq{y}} + \int_{P_t}\tenq{y}\wedge\left[\tenq{b}-\rho\bar{\tenq{a}}\right]d\volume_{\tenq{y}} = \tenq{0}\]

by divergence theorem

\[\Rightarrow \int_{P_t}\left\{\nabla\cdot\tenq{B}+\tenq{y}\wedge\left[\tenq{b}-\rho\bar{\tenq{a}}\right]\right\}d\volume_{\tenq{y}} = \tenq{0}\]

localization argument

\[\Rightarrow \nabla\cdot\tenq{B}+\tenq{y}\wedge\left[\tenq{b}-\rho\bar{\tenq{a}}\right] = \tenq{0}\]

but 
\[\left[\nabla\cdot\tenq{B}\right]_i = \frac{\partial B_{il}}{\partial y_l} = \epsilon_{ijk}\frac{\partial \tau_{kl}}{\partial y_l}y_j + \epsilon_{ijk}\tau_{kl}\delta_{jl}\]

\[\Rightarrow \nabla\cdot\tenq{B} = \tenq{y}\wedge\left(\nabla\cdot\tenq{\tau}\right) + \epsilon_{ijk}\tau_{kj}\tenq{e}_i\]

\[\therefore \epsilon_{ijk}\tau_{kj}\tenq{e}_i+\tenq{y}\wedge\cancelto{0 \text{ from }\eqref{Local version of balance laws-*}}{\left[\nabla\cdot\tenq{\tau}+\tenq{b}-\rho\bar{\tenq{a}}\right]} = \tenq{0}\]

\[\Rightarrow \epsilon_{ijk}\tau_{kj}\tenq{e}_i = \tenq{0}\text{ or }\epsilon_{ijk}\tau_{kj} = 0\]
\[\Rightarrow \tenq{\tau} = \tenq{\tau}^T\cdots\text{symmetric tensor}\]

\end{proof}

However, if we analyze ferrofluid, $\tenq{\tau} \neq \tenq{\tau}^T$ because of the effect of \emph{body moment density} $\tenq{m}$. The balance of angular momentum would be rewritten as 

\[\int_{P_t}\tenq{m}d\volume_{\tenq{y}} + \int_{\partial P_t}\tenq{y}\wedge\tenq{\tau}\tenq{n}dA_{\tenq{y}} + \int_{P_t}\tenq{y}\wedge\left[\tenq{b}-\rho\bar{\tenq{a}}\right]d\volume_{\tenq{y}} = \tenq{0}\]

\begin{remark}\mbox{}\\

\begin{enumerate}[label=(\arabic*)]

\item
equation of motion (local version of conservation of linear momentum)
\[\nabla\cdot\tenq{\tau} + \tenq{b} = \rho\underbrace{\left[\frac{\partial \bar{\tenq{v}}}{\partial t} + \left(\nabla\bar{\tenq{v}}\right)\bar{\tenq{v}}\right]}_{ = \bar{\tenq{a}}}\]

\item
For \emph{rest motions} where $\tenq{y} = \hat{\tenq{y}}(\tenq{x})$ (independent of $t$) then $\bar{\tenq{v}} = \bar{\tenq{a}} = \tenq{0}$. 

\[\nabla\cdot\tenq{\tau} + \tenq{b} = \tenq{0} \text{ or }\tau_{ij,j} + b_i = 0\cdots\text{equilibrium equation}\]

\item
Since $\tenq{\tau} = \tenq{\tau}^T$ (symmetric)

\[\tenq{\tau} = \sum_{i=1}^3\tau_i\tenq{e}_i\otimes\tenq{e}_i\]

where 
\begin{align*}
\tau_i&\cdots\text{principal values of }\tenq{\tau}\text{ or principal stresses}\\
\tenq{e}_i&\cdots\text{principal (spatial) axes or stresses}
\end{align*}

\item
useful relation:\\
recall the formula $\left(\tenq{u}\otimes\tenq{v}\right)\tenq{w} = \tenq{u}\left(\tenq{v}\cdot\tenq{w}\right)$

\begin{figure}[H]
\centering
\resizebox{0.4\textwidth}{!}{\input{figure32.eps_tex}}
\caption{Scheme of orientation of traction}
\label{Scheme of orientation of traction}
\end{figure}

\[
\left\{
\begin{aligned}
\tenq{t}_s &= \tenq{t}(\tenq{n})-\left(\tenq{t}(\tenq{n})\cdot\tenq{n}\right)\tenq{n}\\
\tenq{t}_n &= \left(\tenq{t}(\tenq{n})\cdot\tenq{n}\right)\tenq{n}
\end{aligned}
\right.
\Rightarrow
\left\{
\begin{aligned}
\tenq{t}_s &= \left(\tenq{1}-\tenq{n}\otimes\tenq{n}\right)\tenq{t}\\
\tenq{t}_n &= \left(\tenq{n}\otimes\tenq{n}\right)\tenq{t} = \left(\tenq{n}\cdot\tenq{\tau}\tenq{n}\right)\tenq{n}
\end{aligned}
\right.
\]

\end{enumerate}

\end{remark}

\begin{theorem}{Power theorem}

Let $\hat{\tenq{y}}(\bullet,t)\colon R\to R_t$ be a motion with spatial velocity $\bar{\tenq{v}}$ and $\tenq{D} = \textit{sym }\nabla\bar{\tenq{v}}$. Then

\[\int_{\partial P_t}\tenq{t}(\tenq{n})\cdot\bar{\tenq{v}}dA_{\tenq{y}} + \int_{P_t}\tenq{b}\cdot\bar{\tenq{v}}d\volume_{\tenq{y}} = \int_{P_t}\tenq{\tau}\cdot\tenq{D} d\volume_{\tenq{y}} + \frac{d}{dt}\int_{P_t}\frac{1}{2}\rho\abs{\bar{\tenq{v}}}^2 d\volume_{\tenq{y}}\]

\end{theorem}

\begin{proof}\mbox{}\\

\begin{equation}
\tenq{t}(\tenq{n})\cdot\bar{\tenq{v}} = \left(\tenq{\tau}\tenq{n}\right)\cdot\bar{\tenq{v}} = \tenq{n}\cdot\tenq{\tau}^T\bar{\tenq{v}} = \left(\tenq{\tau}\bar{\tenq{v}}\right)\cdot\tenq{n}
\tag{1}\label{Power theorem-(1)}
\end{equation}

L.H.S. of identity to be proven. 

\[
r(t) \vc{=}{\eqref{Power theorem-(1)}}{\text{divergence theorem}} \int_{P_t}\left\{\nabla\cdot\left(\tenq{\tau}\bar{\tenq{v}}\right) + \tenq{b}\cdot\bar{\tenq{v}}\right\}d\volume_{\tenq{y}}
\]

\[\nabla\cdot\left(\tenq{\tau}\bar{\tenq{v}}\right) = \left(\tau_{ij}\bar{v}_{j}\right)_{,i} = \tau_{ij,i}\bar{v}_j + \tau_{ij}\bar{v}_{j,i} = \tau_{ji,i}\bar{v}_j + \tau_{ij}\bar{v}_{j,i}\]

\begin{align*}
\nabla\cdot\left(\tenq{\tau}\bar{\tenq{v}}\right)
&= \left(\nabla\cdot\tenq{\tau}\right)\cdot\bar{\tenq{v}} + \tenq{\tau}\cdot\nabla\bar{\tenq{v}}\\
&= \left(\nabla\cdot\tenq{\tau}\right)\cdot\bar{\tenq{v}} + \tenq{\tau}\cdot\left[\tenq{D} + \tenq{W}\right]\\
&= \left[\nabla\cdot\tenq{\tau}\right]\cdot\bar{\tenq{v}} + \tenq{\tau}\cdot\tenq{D}
\end{align*}

where

\begin{align*}
\tenq{D}&= \textit{sym }\nabla\bar{\tenq{v}}\\
\tenq{W}&= \textit{skew }\nabla\bar{\tenq{v}}
\end{align*}

\begin{align*}
\Rightarrow
r(t)
&= \int_{P_t}\left\{\underbrace{\left[\nabla\cdot\tenq{\tau} + \tenq{b}\right]}_{=\rho\bar{\tenq{a}}}\cdot\bar{\tenq{v}} + \tenq{\tau}\cdot\tenq{D}\right\}d\volume_{\tenq{y}}\\
&= \int_{P_t}\left\{\rho\bar{\tenq{a}}\cdot\bar{\tenq{v}} + \tenq{\tau}\cdot\tenq{D}\right\}d\volume_{\tenq{y}}
\end{align*}

recall the formula $\frac{1}{2}\frac{D}{Dt}\left(\bar{\tenq{v}}\cdot\bar{\tenq{v}}\right) = \frac{1}{2}\frac{D}{Dt}\abs{\bar{\tenq{v}}}^2$

\[
\bar{\tenq{a}}\cdot\bar{\tenq{v}} = \frac{1}{2}\vc{\left(\bar{\tenq{v}}\cdot\bar{\tenq{v}}\right)}{\bullet}{} = \vc{\left\{\frac{1}{2}\abs{\bar{\tenq{v}}}^2\right\}}{\bullet}{}
\]

\begin{align*}
r(t)
&= \int_{P_t}\tenq{\tau}\cdot\tenq{D}d\volume_{\tenq{y}} + \int_{P_t}\rho\vc{\left\{\frac{1}{2}\abs{\bar{\tenq{v}}}^2\right\}}{\bullet}{}d\volume_{\tenq{y}}\\
&\vc{=}{RTT}{}\int_{P_t}\tenq{\tau}\cdot\tenq{D}d\volume_{\tenq{y}} + \frac{d}{dt}\int_{P_t}\frac{1}{2}\rho\abs{\bar{\tenq{v}}}^2d\volume_{\tenq{y}}
\end{align*}
\end{proof}

\begin{remark}\mbox{}\\

\begin{align*}
\int_{\partial P_t}\tenq{t}\cdot\bar{\tenq{v}}dA_{\tenq{y}}
&\cdots\text{power (rate of work) done by traction on }\partial P_t\\
\int_{P_t}\tenq{b}\cdot\bar{\tenq{v}}d\volume_{\tenq{y}}
&\cdots\text{power done by body force}\\
\int_{P_t}\frac{1}{2}\rho\abs{\bar{\tenq{v}}}^2d\volume_{\tenq{y}}
&\cdots\text{kinetic energy}\\
\int_{P_t}\tenq{\tau}\cdot\tenq{D} d\volume_{\tenq{y}}
&\cdots\text{stress power}
\end{align*}

\end{remark}